\section{实战案例：如何判断 ``+0.6 分'' 到底算不算进步}

\subsection{案例设定}
为了把课程中的统计思想落到工程决策，我们构造一个典型场景：团队有两个版本，Model-A（基线）和 Model-B（候选），在同一 benchmark、同一协议下各跑 10 次。平均分看起来 Model-B 高 0.6 分，但 run-to-run 波动也明显。

\begin{table}[H]
\centering
\begin{tabular}{p{2.5cm}p{2.2cm}p{2.2cm}p{2.0cm}p{2.2cm}}
\toprule
\textbf{模型} & \textbf{Mean} & \textbf{Std} & \textbf{Runs} & \textbf{95\% CI} \\
\midrule
Model-A & 72.4 & 1.1 & 10 & [71.7, 73.1] \\
Model-B & 73.0 & 1.2 & 10 & [72.2, 73.8] \\
\bottomrule
\end{tabular}
\caption{单看均值似乎 Model-B 更强，但置信区间仍有重叠}
\end{table}

\begin{knowledgebox}{课程强调的决策规则}
当两个模型置信区间重叠明显时，不应直接下结论 ``B 胜过 A''。更稳妥做法是：增加重复次数、做 paired test、并在关键子集上复核。
\end{knowledgebox}

\subsection{从总体分数转向分桶结论}
把样本按难度分桶后，我们可能看到完全不同的结论：Model-B 在 easy split 提升很大，但在 hard split 几乎无提升。若真实业务更接近 hard split，那么总体 +0.6 分并不代表真实价值。

\begin{table}[H]
\centering
\begin{tabular}{p{2.3cm}p{2.4cm}p{2.4cm}p{2.4cm}}
\toprule
\textbf{难度分桶} & \textbf{Model-A} & \textbf{Model-B} & \textbf{差值 (B-A)} \\
\midrule
Easy & 84.1 & 85.7 & +1.6 \\
Medium & 71.5 & 72.0 & +0.5 \\
Hard & 58.8 & 58.9 & +0.1 \\
\bottomrule
\end{tabular}
\caption{分桶后可见改进主要来自 easy split，而非 hard split}
\end{table}

\begin{warningbox}{只报总体分数会导致错误上线}
如果业务请求以 hard case 为主，盲目按总体分数上线，可能带来 ``线上体感无提升甚至退化'' 的结果。课程中把这类问题称为 benchmark interpretation failure。
\end{warningbox}

\subsection{把显著性检验写进发布流程}
在团队流程上，建议把 ``统计显著性门槛'' 变成 release gate。候选模型只有在关键任务集上满足预设门槛（例如 p-value 与 effect size 同时达标）才可进入灰度。这让评测从 ``结果展示'' 升级为 ``质量控制''。

\begin{importantbox}{可执行的发布门槛示例}
\begin{itemize}
    \item 关键任务集上 $\Delta \text{score} \ge 0.8$；
    \item paired bootstrap 的 95\% CI 下界 $>0$；
    \item hard split 不得退化超过 0.2 分；
    \item 工具调用效率指标（tokens / tool calls / latency）不得恶化超过阈值。
\end{itemize}
\end{importantbox}

\subsection{本章小结}
``+0.6 分'' 的意义取决于方差、分桶和业务分布。统计显著性不是学术装饰，而是模型发布的风险控制工具。

